\section{Working environment}
\label{section:preparation:environment}
  
The setup of a working environment can have a big impact on a software's performance. Clearly the hardware has a large impact, but so can the software environment not to mention the structure of the cluster, e.g., the wait times on partitions can impact the turnaround time for the assessment.

In this section we lay out some features of the working environment that we believe are important when setting up the working environment for code compilation and benchmarking, starting with the machine. The hardware details are crucially important as we may later on wish to do comparative benchmarking on different hardware, or the code submitter may have requested a hardware platform that they have access to or are interested in.

\subsection{Machine details}
\label{subsection:preparation:machine}

Modern supercomputers are very complex and it is virtually impossible to
memorise all the hardware details of a testbed given the zoo of processor types,
generations and flavours.
While it is possible to consult your compute centre's webpages or to consult
your system administrator, it might be more reasonable to obtain that data
directly from the machine.

\begin{assessmentcrime}
 Never grab the specification from the login node. Instead, ask for an
 interactive terminal on a compute node and get the data there. Many
 supercomputers' login nodes are slightly different to the compute nodes,
 i.e.~slightly different make, more memory, more cores, \ldots So you might end
 up with the wrong hardware specification if you analyse a login node.
\end{assessmentcrime}

\begin{instructions}{\linux}
 On every Linux system, information about the CPU on the system is held in a
 central file.
 You can read it out by typing \texttt{cat /proc/cpuinfo} into the terminal.
\end{instructions}


\begin{instructions}{\likwid}
 If you have Likwid installed, you can invoke \texttt{likwid-topology} to obtain
 in-depth information about your CPU.
\end{instructions}

\begin{instructions}{\maqao}
 In \maqao, you find some topology information in the menu rubric \emph{Topology}. 
\end{instructions}

In the checklist below we recommend some hardware features to note:
\begin{checklist}
\item Number of nodes
\item The node layout, e.g., dual socket of AMD EPYC 7702 64-Core Processors
\item Caches (including sizes and associated cores) and memory 
\item NUMA domains
\end{checklist}

\paragraph*{How it works}

\begin{itemize}[noitemsep]
  \item[$\boxplus$] Gather the exact spec of your testbed under realistic load. If this information is already available (e.g.~due to previous assessment exercises), it makes sense to reuse it.
  \item[$\boxplus$] Note down any features of the testbeds queues and partitions, including which are most likely to be used for the assessment.
\end{itemize}


\paragraph*{How it does not work}

\begin{itemize}[noitemsep]
  \item[$\boxminus$] Throw advanced tools onto your code and expect that they report the theoretical machine capabilities themselves.
\end{itemize}

\subsection{Software environment}
\label{subsection:preparation:software-environment}

The software environment, e.g., compilers, libraries, modules, etc. can have a significant impact on a code's performance. We do not yet worry about the impact that these dependencies have on the code's performance, as by this stage we are simply interested in building the code. However, as with machine details, it is very useful to have the software environment documented.

\begin{instructions}{Modules}
To grab the loaded modules and dependencies (following the code's documentation) we can simply 
\begin{verbatim}
module list
\end{verbatim}
\end{instructions}

\paragraph*{How it works}

\begin{itemize}[noitemsep]
  \item[$\boxplus$] Following the documentation, load in relevant dependencies and note them down using the \texttt{module list} command including any version numbers.
\end{itemize}


\paragraph*{How it does not work}

\begin{itemize}[noitemsep]
  \item[$\boxminus$] Copy across the dependencies as they appear in the documentation.
\end{itemize}


